\documentclass{beamer}

\usepackage[spanish]{babel}
\usepackage[utf8]{inputenc}
\usepackage{amsmath}
\usetheme{metropolis}% modern fork of the metropolis theme
\usepackage{booktabs}
\usepackage{float}
\usepackage{amsmath}
\usepackage{amsfonts}


\definecolor{mPurple}{RGB}{91,48,105}
\setbeamercolor{palette primary}{bg=mPurple}
\setbeamercolor{normal text}{fg=black}
%\usetheme{Madrid}

\title{Stress Testing en Sistemas de Pensiones}
\subtitle{Módulo 2: Identificación y Medición de Riesgos Previsionales}
\author{MACROPOL}
\date{\today}

\begin{document}

%------------------------------------------------
\frame{\titlepage}
%\begin{frame}{Contenido}
%\tableofcontents    
%\end{frame}
%------------------------------------------------
\section{Módulo 2: Identificación y Medición de Riesgos Previsionales}

%------------------------------------------------
\begin{frame}[t]{Objetivo del Módulo}
\begin{itemize}

\item Identificar los principales riesgos que afectan los sistemas de pensiones
\item Medir su impacto sobre variables previsionales clave
\item Construir la base para escenarios de stress testing
\end{itemize}

\vspace{0.3cm}
\textbf{Idea central:} No se puede evaluar lo que no se mide correctamente
\end{frame}
%-----------------------------------------------
\begin{frame}{Concepto de riesgo en pensiones y Stress Testing}
\begin{itemize}

    \item En los sistemas de pensiones de capitalización individual, el concepto de riesgo es central porque determina la sostenibilidad de pensiones futuras y la estabilidad del sistema.
    \vspace{0.3cm}
    \item Por \textbf{riesgo previsional} debe entenderse como una combinación de incertidumbre económica, financiera y demográfica, pero tratada de forma cuantificable mediante herramientas de stress testing.
\end{itemize}   
\end{frame}
%-----------------------------------------------

\begin{frame}{Riesgo vs. Incertidumbre}

\begin{itemize}

      \item \textbf{Riesgo:} situaciones donde los resultados posibles y sus probabilidades pueden ser estimados $\rightarrow$ \textbf{medible y modelable}
      \vspace{0.2cm}
      \item \textbf{Incertidumbre:} situaciones donde no se conocen las probabilidades $\rightarrow$ \textbf{no modelable directamente}
      \vspace{0.2cm}
      \item Intuición económica:      
         \begin{itemize}
            \item En pensiones operan bajo horizonte largo
            \vspace{0.2cm}
            \item riesgo = volatilidad de retornos, longevidad estimable
            \vspace{0.2cm}
            \item incertidumbre = cambios estructurales (reformas, shocks extremos)
        \end{itemize}
\end{itemize}
  
\end{frame}

%-------------------------------------------------
\begin{frame}{Riesgo en finanzas}

    \begin{itemize}
        \item El riesgo se define como variabilidad de los retornos esperados
        \vspace{0.3cm}
        \item Medido por: 
        \begin{itemize}
            \item varianza, 
            \item volatilidad
            \item VaR
        \end{itemize}
        \vspace{0.3cm}
        \item Supuesto: estabilidad estadística    
        \vspace{0.3cm}
        \item Limitaciones del enfoque tradicional
        \begin{itemize}
            \item Subestima eventos extremos
            \item No captura crisis sistémicas
            \item Dependencia de datos históricos
        \end{itemize}
\end{itemize} 
    
\end{frame}
%------------------------------------------------
\begin{frame}{Riesgos Previsionales: Definición y Alcance}
\footnotesize
\begin{itemize}

\item \textbf{Riesgo de mercado}: Pérdidas potenciales en el valor del portafolio previsional debido a fluctuaciones en precios de activos financieros (acciones, bonos, tipo de cambio). Afecta directamente la acumulación del fondo individual.

\vspace{0.3cm}

\item \textbf{Riesgo de longevidad}: 
Riesgo de que los afiliados vivan más de lo esperado, aumentando la duración del pago de pensiones y reduciendo la tasa de reemplazo esperada.

\vspace{0.3cm}

\item \textbf{Riesgo de tasa de interés}: 
Impacto de cambios en la estructura temporal de tasas sobre el valor de los activos y las rentas vitalicias, afectando tanto acumulación como desacumulación.

\vspace{0.3cm}

\item \textbf{Riesgo de densidad de cotización}:
Incertidumbre asociada a la frecuencia y continuidad de las contribuciones. Bajos niveles de cotización reducen el saldo acumulado final.

\vspace{0.3cm}

\item \textbf{Riesgo regulatorio}:
Cambios en normativa (comisiones, edad de retiro, tasas de contribución, reglas de inversión) que afectan los retornos esperados y la sostenibilidad del sistema.

\end{itemize}
\end{frame}

%-----------------------------------------------
\begin{frame}[t]{Característica del riesgo previsional}
\vspace{0.5cm}
\begin{itemize}
    \item Horizonte de largo plazo
    \item Interacción macro-financiera-demográfica
    \item Impacto en tasa de reemplazo
\end{itemize}
    
\end{frame}
%------------------------------------------------

\begin{frame}[t]{Riesgo diversificable vs. no diversificable}

\textbf{Distinción clave en teoría de portafolio y stress testing}

\begin{itemize}

    \item \textbf{Riesgo diversificable (idiosincrático)}  
    \begin{itemize}
        \item Riesgo específico de un activo, emisor o sector
        \item No está correlacionado perfectamente con el mercado
        \item \textbf{Definición:} componente del riesgo que puede eliminarse mediante diversificación
        \item $\rightarrow$ \text{Se reduce al aumentar el número de activos en el portafolio}
    \end{itemize}
    
   
    \item \textbf{Riesgo no diversificable (sistémico)}  
    \begin{itemize}
        \item Riesgo común a todos los activos (factores macroeconómicos)
        \item Ejemplos: shocks de tasas de interés, inflación, crisis financieras
        \item \textbf{Definición:} componente del riesgo que no puede eliminarse mediante diversificación
        \item $\rightarrow$ \text{Determinado por factores agregados del sistema}
    \end{itemize}
     
\end{itemize}

\end{frame}
%----------------------------------------------------
\begin{frame}[t]{Riesgo diversificable vs. no diversificable}

  \textbf{Aplicación en pensiones}
    \begin{itemize}
        \item Portafolios previsionales altamente expuestos a \textbf{riesgo sistémico global}
        \item \textbf{Mercado financiero:} principal fuente de riesgo no diversificable
        \item \textbf{Riesgo de longevidad:} 
        \begin{itemize}
            \item componente idiosincrático (individual)
            \item componente sistémico (mejoras estructurales en esperanza de vida)
        \end{itemize}
    \end{itemize}

\vspace{0.3cm}
\textit{Implicación para stress testing: los escenarios adversos deben modelar principalmente shocks sobre factores sistémicos.}
    
\end{frame}
%----------------------------------------------------
\begin{frame}{Riesgo en Stress Testing}
\begin{itemize}
    \item El stress testing redefine el riesgo como:
    riesgo = pérdida bajo escenarios extremos plausibles
    \vspace{0.3cm}
    \item No depende solo de probabilidades históricas
    \vspace{0.3cm}
    
    \item Se enfoca en:
    \begin{itemize}
        \item eventos raros
        \item colas de distribución
        \item interacciones macro-financieras
    \end{itemize}
    \item Rol clave en regulación
    \begin{itemize}
        \item evaluar resiliencia del sistema
        \item identificar vulnerabilidades sistémicas
    \end{itemize}
\end{itemize}
    
\end{frame}
% ------------------------------------------------
\begin{frame}{Diferencias clave}
\begin{itemize}
    \item Finanzas tradicionales: basado en probabilidades
    \item Stress testing: basado en escenarios, enfoque forward looking
\end{itemize}
    
\end{frame}

%------------------------------------------------
\begin{frame}{Canales de transmisión del riesgo previsional}
\textbf{Mecanismos de impacto}

\begin{itemize}
\item \textbf{Portafolio} $\rightarrow$ valor del fondo acumulado
\item \textbf{Mercado laboral} $\rightarrow$ densidad de cotización
\item \textbf{Demografía} $\rightarrow$ duración de beneficios
\item \textbf{Tasas de interés} $\rightarrow$ valorización de activos y rentas
\item \textbf{Regulación} $\rightarrow$ reglas de contribución y retiro
\end{itemize}

\vspace{0.3cm}
\textbf{Variable final de interés:}
\begin{itemize}
\item Tasa de reemplazo (resultado integrado del sistema)
\end{itemize}

\vspace{0.2cm}
\textbf{Clave:} interacción simultánea de riesgos
\end{frame}

%------------------------------------------------
\begin{frame}{Riesgo de Mercado}
\textbf{Intuición económica}
\begin{itemize}
\item Caídas en precios $\rightarrow$ menor acumulación de ahorro
\item Alta correlación en crisis $\rightarrow$ pérdida sistémica
\end{itemize}

\vspace{0.3cm}
\textbf{Modelo cuantitativo}
\begin{itemize}
\item Retornos logarítmicos: $r_t = \ln(P_t/P_{t-1})$
\item Dinámica de volatilidad (GARCH)
\item Correlaciones dependientes del ciclo
\end{itemize}

\vspace{0.3cm}
\textbf{Métricas}
\begin{itemize}
\item Volatilidad del portafolio
\item Value at Risk (VaR)
\item Pérdidas bajo stress (escenarios)
\end{itemize}
\end{frame}

%------------------------------------------------
\begin{frame}{Riesgo de Longevidad}
\textbf{Intuición económica}
\begin{itemize}
\item Mayor esperanza de vida $\rightarrow$ mayor período de pago de pensiones
\item Riesgo sistémico (no diversificable completamente)
\end{itemize}

\vspace{0.3cm}
\textbf{Modelo actuarial}
\begin{itemize}
\item Tablas de mortalidad
\item Modelo Lee-Carter:
\[
\ln m_{x,t} = a_x + b_x k_t + \epsilon_{x,t}
\]
\end{itemize}

\vspace{0.3cm}
\textbf{Impacto}
\begin{itemize}
\item Reducción de tasa de reemplazo
\item Incremento del pasivo previsional implícito
\end{itemize}
\end{frame}

%------------------------------------------------
\begin{frame}{Riesgo de Tasa de Interés}
\textbf{Intuición económica}
\begin{itemize}
\item Cambios en tasas afectan:
\begin{itemize}
\item Valorización de bonos
\item Tasas de descuento de pensiones
\end{itemize}
\end{itemize}

\vspace{0.3cm}
\textbf{Modelo}
\begin{itemize}
\item Estructura temporal de tasas (curva yield)
\item Modelos: Vasicek / CIR
\end{itemize}

\vspace{0.3cm}
\textbf{Métrica}
\begin{itemize}
\item Duración y convexidad del portafolio
\item Sensibilidad del valor presente
\end{itemize}
\end{frame}

%------------------------------------------------
\begin{frame}{Riesgo de Densidad de Cotización}
\textbf{Intuición económica}
\begin{itemize}
\item Trayectorias laborales incompletas reducen ahorro acumulado
\item Alta informalidad $\rightarrow$ menor cobertura previsional
\end{itemize}

\vspace{0.3cm}
\textbf{Modelo}
\begin{itemize}
\item Probabilidad de cotizar en el tiempo
\item Procesos estocásticos de empleo/desempleo
\end{itemize}

\vspace{0.3cm}
\textbf{Impacto}
\begin{itemize}
\item Menor saldo acumulado
\item Alta dispersión en tasas de reemplazo
\end{itemize}
\end{frame}

%------------------------------------------------
\begin{frame}{Riesgo Macroeconómico}
\textbf{Intuición económica}
\begin{itemize}
\item Choques macro afectan simultáneamente:
\begin{itemize}
\item Retornos financieros
\item Empleo y salarios
\end{itemize}
\end{itemize}

\vspace{0.3cm}
\textbf{Variables clave}
\begin{itemize}
\item PIB
\item Inflación
\item Tasas de interés
\end{itemize}

\vspace{0.3cm}
\textbf{Herramienta}
\begin{itemize}
\item Modelos VAR / escenarios macro-financieros
\end{itemize}
\end{frame}

%------------------------------------------------
\begin{frame}{Riesgo de Liquidez}
\textbf{Intuición económica}
\begin{itemize}
\item Necesidad de liquidar activos en condiciones adversas
\end{itemize}

\vspace{0.3cm}
\textbf{Escenarios relevantes}
\begin{itemize}
\item Retiros masivos
\item Crisis financieras
\end{itemize}

\vspace{0.3cm}
\textbf{Métrica}
\begin{itemize}
\item Ratio activos líquidos / obligaciones
\item Haircuts en escenarios de stress
\end{itemize}
\end{frame}

%------------------------------------------------
\begin{frame}{Riesgo Regulatorio}
\textbf{Intuición económica}
\begin{itemize}
\item Cambios en reglas alteran incentivos y resultados
\end{itemize}

\vspace{0.3cm}
\textbf{Ejemplos}
\begin{itemize}
\item Cambios en tasas de contribución
\item Retiros extraordinarios
\item Modificación de edad de jubilación
\end{itemize}

\vspace{0.3cm}
\textbf{Naturaleza}
\begin{itemize}
\item No probabilístico (incertidumbre)
\item Se modela vía escenarios
\end{itemize}
\end{frame}

%------------------------------------------------
\begin{frame}{Integración de Riesgos}
\textbf{Clave conceptual}

\begin{itemize}
\item Los riesgos no son independientes
\item Existen efectos de amplificación:
\begin{itemize}
\item Crisis $\rightarrow$ mercado + empleo + liquidez
\end{itemize}
\end{itemize}

\vspace{0.3cm}
\textbf{Enfoque de stress testing}
\begin{itemize}
\item Escenarios conjuntos (macro + financiero + demográfico)
\item Evaluación de impacto en tasa de reemplazo
\end{itemize}

\vspace{0.3cm}
\textbf{Objetivo}
\begin{itemize}
\item Medir resiliencia del sistema previsional
\end{itemize}
\end{frame}


%------------------------------------------------
\begin{frame}{Indicadores de Vulnerabilidad}
\textbf{De riesgo a impacto previsional}

\begin{itemize}
\item \textbf{Métricas financieras}
\begin{itemize}
\item Value at Risk (VaR)
\item Expected Shortfall (ES)
\end{itemize}

\vspace{0.2cm}

\item \textbf{Métricas previsionales}
\begin{itemize}
\item Tasa de reemplazo proyectada
\item Probabilidad de insuficiencia previsional
\end{itemize}
\end{itemize}

\vspace{0.3cm}

\textbf{Clave conceptual:}
\begin{itemize}
\item Shock $\rightarrow$ portafolio $\rightarrow$ saldo $\rightarrow$ pensión $\rightarrow$ RR
\end{itemize}

\vspace{0.3cm}

\textbf{Objetivo:} traducir riesgo en métricas de política pública
\end{frame}

%------------------------------------------------
\begin{frame}{Value at Risk (VaR)}
\textbf{Definición}

\[
VaR_\alpha = \inf \{x : P(L > x) \leq 1-\alpha \}
\]

\vspace{0.3cm}

\textbf{Intuición}
\begin{itemize}
\item Pérdida máxima esperada bajo condiciones normales
\item No captura eventos extremos (colas)
\end{itemize}

\vspace{0.3cm}

\textbf{Limitación en pensiones}
\begin{itemize}
\item Horizonte corto
\item No refleja impacto en tasa de reemplazo
\end{itemize}
\end{frame}

%------------------------------------------------
\begin{frame}{Expected Shortfall (ES)}
\textbf{Definición}

\[
ES_\alpha = E[L \mid L > VaR_\alpha]
\]

\vspace{0.3cm}

\textbf{Intuición}
\begin{itemize}
\item Pérdida promedio en escenarios extremos
\item Captura riesgo de cola
\end{itemize}

\vspace{0.3cm}

\textbf{Ventaja}
\begin{itemize}
\item Métrica coherente de riesgo
\item Más relevante para stress testing
\end{itemize}
\end{frame}

%------------------------------------------------
\begin{frame}{Tasa de Reemplazo (RR)}
\textbf{Definición}

\[
RR = \frac{\text{Pensión}}{\text{Salario}}
\]

\vspace{0.3cm}

\textbf{Determinantes}
\begin{itemize}
\item Retornos del portafolio
\item Densidad de cotización
\item Esperanza de vida
\item Tasas de interés
\end{itemize}

\vspace{0.3cm}

\textbf{Interpretación}
\begin{itemize}
\item Métrica central de suficiencia previsional
\item Variable objetivo del stress testing
\end{itemize}
\end{frame}

%------------------------------------------------
\begin{frame}{Probabilidad de Insuficiencia Previsional}
\textbf{Definición}

\[
P(RR < RR^*)
\]

\vspace{0.3cm}

\textbf{Interpretación}
\begin{itemize}
\item Probabilidad de no alcanzar un umbral mínimo de pensión
\item Métrica clave para regulación y política social
\end{itemize}

\vspace{0.3cm}

\textbf{Enfoque}
\begin{itemize}
\item Requiere simulación (Monte Carlo)
\item Integra múltiples riesgos simultáneamente
\end{itemize}
\end{frame}

%------------------------------------------------
\begin{frame}{De Riesgo Financiero a Riesgo Previsional}
\textbf{Cadena de transmisión}

\begin{itemize}
\item Shock financiero $\rightarrow$ retornos
\item Retornos $\rightarrow$ saldo acumulado
\item Saldo $\rightarrow$ pensión
\item Pensión $\rightarrow$ tasa de reemplazo
\end{itemize}

\vspace{0.3cm}

\textbf{Insight clave}
\begin{itemize}
\item El riesgo relevante no es la pérdida del portafolio
\item Es la caída en la tasa de reemplazo
\end{itemize}
\end{frame}

%------------------------------------------------
\begin{frame}{Enfoque de Stress Testing}
\textbf{Metodología}

\begin{itemize}
\item Definir escenario (baseline vs adverso)
\item Simular trayectorias de variables
\item Evaluar distribución de RR
\end{itemize}

\vspace{0.3cm}

\textbf{Outputs}
\begin{itemize}
\item RR promedio
\item Percentiles (p5, p50, p95)
\item Probabilidad de insuficiencia
\end{itemize}

\vspace{0.3cm}

\textbf{Uso}
\begin{itemize}
\item Evaluación de resiliencia del sistema
\end{itemize}
\end{frame}

%------------------------------------------------
\begin{frame}{Mensaje Clave del Módulo}
\begin{itemize}
\item El riesgo financiero es solo un medio
\item El objetivo es medir impacto previsional
\item Stress testing conecta ambos mundos
\end{itemize}

\vspace{0.3cm}

\textbf{Resultado final}
\begin{itemize}
\item Evaluación de suficiencia y sostenibilidad
\end{itemize}
\end{frame}

%------------------------------------------------
\begin{frame}{LAB 1: Simulación de Riesgo Previsional}

\textbf{Objetivo}
\vspace{0.3cm}
Cuantificar el impacto del riesgo financiero sobre:

\begin{itemize}
\item Saldo acumulado
\item Pensión
\item Tasa de reemplazo (RR)
\end{itemize}

\vspace{0.3cm}

\textbf{Enfoque}
\begin{itemize}
\item Simulación Monte Carlo
\item Baseline vs escenario adverso
\end{itemize}
\end{frame}

%------------------------------------------------
\begin{frame}{Modelo Conceptual}

\textbf{Canal de transmisión}

\[
\text{Retornos} \rightarrow \text{Saldo} \rightarrow \text{Pensión} \rightarrow RR
\]

\vspace{0.5cm}

\textbf{Ecuaciones clave}

\begin{itemize}
\item Acumulación:
\[
W_{t+1} = W_t (1 + r_t) + c \cdot salario
\]

\item Pensión:
\[
P = \frac{W_T}{a}
\]

\item Tasa de reemplazo:
\[
RR = \frac{P}{salario}
\]
\end{itemize}
\end{frame}

%------------------------------------------------
\begin{frame}{Supuestos del Modelo (Baseline)}

\begin{itemize}
\item Horizonte: 30 años
\item Salario constante
\item Tasa de contribución: 10\%
\item Retornos: normales
\item Factor de anualidad fijo
\end{itemize}

\vspace{0.3cm}

\textbf{Parámetros}

\begin{itemize}
\item $\mu = 5\%$
\item $\sigma = 12\%$
\item $a = 15$
\end{itemize}

\end{frame}

%------------------------------------------------
\begin{frame}[fragile]{Implementación en R (Simulación)}
\scriptsize
    \begin{verbatim}
    set.seed(123)
    
    T <- 30
    n_sim <- 5000
    
    salario <- 1000
    contribucion <- 0.10 * salario
    mu <- 0.05
    sigma <- 0.12
    a <- 15
    
    RR <- numeric(n_sim)
    
    for(i in 1:n_sim){
      W <- 0
      
      for(t in 1:T){
        r <- rnorm(1, mu, sigma)
        W <- W * (1 + r) + contribucion
      }
      
      pension <- W / a
      RR[i] <- pension / salario
    }
    \end{verbatim}
\end{frame}

%------------------------------------------------
\begin{frame}[fragile]{Resultados}
\textbf{Cálculo de indicadores}

\vspace{0.3cm}

\textbf{En R:}
\scriptsize
\begin{verbatim}
mean(RR)
mean(RR < 0.4)
\end{verbatim}

\vspace{0.4cm}

\textbf{Interpretación}

\begin{itemize}
\item RR promedio
\item Probabilidad de insuficiencia: $P(RR < 0.4)$
\end{itemize}
\end{frame}

%------------------------------------------------
\begin{frame}{Ejercicio: Escenario Adverso}

\textbf{Modificar parámetros}

\begin{itemize}
\item $\mu = -2\%$
\item $\sigma = 20\%$
\end{itemize}

\vspace{0.3cm}

\textbf{Preguntas}

\begin{itemize}
\item ¿Cómo cambia la distribución de RR?
\item ¿Qué ocurre con $P(RR < 0.4)$?
\item ¿Qué riesgo domina?
\end{itemize}
\end{frame}
%------------------------------------------------
\begin{frame}{Interpretación Económica}
\begin{itemize}
\item Mayor volatilidad $\rightarrow$ mayor dispersión de RR
\item Menores retornos $\rightarrow$ menor pensión
\item Aumenta la probabilidad de insuficiencia
\end{itemize}

\vspace{0.3cm}

\textbf{Insight clave}

\begin{itemize}
\item El riesgo relevante no es la pérdida financiera
\item Es la caída en la tasa de reemplazo
\end{itemize}
\end{frame}

%------------------------------------------------
\begin{frame}{Conexión con Stress Testing}
\begin{itemize}
\item Comparación baseline vs adverso
\item Evaluación de distribuciones
\item Medición de vulnerabilidad previsional
\end{itemize}

\vspace{0.3cm}

\textbf{Output clave}

\[
P(RR < RR^*)
\]

\vspace{0.3cm}

\textbf{Aplicación}

\begin{itemize}
\item Regulación
\item Diseño de políticas previsionales
\end{itemize}
\end{frame}
%----------------------------------------------
\begin{frame}{De Riesgos a Escenarios}

\textbf{Problema:}  
Los riesgos identificados son abstractos

\vspace{0.3cm}

\textbf{Solución:}  
Traducir riesgos en escenarios macro-financieros coherentes

\vspace{0.3cm}

\begin{itemize}
\item Riesgos $\rightarrow$ Factores (retornos, inflación, longevidad)
\item Factores $\rightarrow$ Shocks (magnitud + correlación)
\item Shocks $\rightarrow$ Escenarios (trayectorias en el tiempo)
\end{itemize}

\vspace{0.4cm}

\textbf{Idea clave:}  
Un stress test evalúa estados del mundo, no shocks aislados

\end{frame}

%---------------------------------------
\begin{frame}{Limitaciones del Enfoque Determinístico}

\begin{itemize}
\item Un solo escenario ignora la incertidumbre real
\item Resultados altamente sensibles a supuestos puntuales
\item No captura eventos extremos (riesgo de cola)
\item No permite asignar probabilidades a resultados
\end{itemize}

\vspace{0.4cm}

\textbf{Problema central:}  
No sabemos qué tan probable es el escenario adverso

\end{frame}
%------------------------------
\begin{frame}{Motivación para Simulación}

\begin{itemize}
\item El sistema previsional enfrenta múltiples futuros posibles
\item Los shocks son dinámicos y acumulativos
\item Las correlaciones importan (crisis = shocks conjuntos)
\end{itemize}

\vspace{0.4cm}

\textbf{Necesidad:}  
Pasar de un escenario único a una distribución de resultados

\end{frame}
%----------------------------------
\begin{frame}{Simulación Monte Carlo}

\begin{itemize}
\item Genera múltiples trayectorias consistentes con los datos
\item Incorpora volatilidad, correlaciones y no linealidades
\item Aproxima la distribución completa de resultados
\end{itemize}

\vspace{0.4cm}

\textbf{Aplicación en pensiones:}
\begin{itemize}
\item Evolución del fondo acumulado
\item Tasa de reemplazo
\item Probabilidad de insuficiencia
\end{itemize}

\end{frame}
%---------------------------------
\begin{frame}{Modelo Conceptual del Sistema Previsional}

\[
W_{t+1}^{(i)} = W_t^{(i)} (1 + r_t^{(i)}) + c_t
\]

\begin{itemize}
\item Cada $i$ representa un escenario posible
\item $r_t$ incorpora shocks macro-financieros
\end{itemize}

\vspace{0.4cm}

\textbf{Resultado:}
\begin{itemize}
\item Distribución de riqueza final $W_T$
\item Distribución de tasa de reemplazo
\end{itemize}

\end{frame}
%----------------------------------
\begin{frame}{De Escenarios a Distribuciones}

\begin{itemize}
\item Cada simulación genera un resultado distinto
\item El conjunto define una distribución de resultados
\end{itemize}

\vspace{0.4cm}

\textbf{Métricas clave:}
\begin{itemize}
\item Valor esperado
\item VaR (riesgo cuantílico)
\item Expected Shortfall (riesgo extremo)
\item Probabilidad de insuficiencia
\end{itemize}

\end{frame}
%------------------------------
\begin{frame}{Mensaje Clave}

\begin{itemize}
\item El riesgo previsional es multidimensional
\item La incertidumbre es inherente al sistema
\end{itemize}

\vspace{0.5cm}

\centering
\textbf{Un stress test transforma escenarios en decisiones de política}

\end{frame}
%--------------------------
%------------------------------------------------
\section{Ejercicio Integrador: De Riesgos a Escenarios}

%------------------------------------------------
\begin{frame}{Ejercicio: Construcción de Escenarios}

\textbf{Objetivo:}  
Traducir riesgos previsionales en escenarios macro-financieros consistentes

\vspace{0.4cm}

\textbf{Contexto:}
\begin{itemize}
\item Una AFP desea evaluar su vulnerabilidad
\item Se han identificado los riesgos principales:
\begin{itemize}
\item Riesgo de mercado
\item Inflación
\item Longevidad
\end{itemize}
\end{itemize}

\vspace{0.4cm}

\textbf{Tarea:}  
Construir escenarios y evaluar su impacto en la tasa de reemplazo

\end{frame}

%------------------------------------------------
\begin{frame}{Estructura del Ejercicio}

\begin{enumerate}
\item Definir factores de riesgo
\item Construir escenarios (baseline vs adverso)
\item Traducir narrativa a shocks
\item Simular trayectorias
\item Analizar resultados
\item Extensión: simulación Monte Carlo
\end{enumerate}

\vspace{0.4cm}

\textbf{Resultado esperado:}  
Comprender cómo los escenarios afectan la suficiencia previsional

\end{frame}

%------------------------------------------------
\begin{frame}{Parte 1: De Riesgos a Factores}

\textbf{Pregunta 1:}  
¿Cómo se traduce cada riesgo en una variable observable?

\vspace{0.4cm}

\begin{itemize}
\item Mercado $\rightarrow$ retorno del portafolio ($r_t$)
\item Inflación $\rightarrow$ inflación ($\pi_t$)
\item Longevidad $\rightarrow$ años de retiro ($L$)
\end{itemize}

\vspace{0.4cm}

\textbf{Pregunta 2:}  
¿Por qué es necesario pasar de riesgos abstractos a variables medibles?

\end{frame}

%------------------------------------------------
\begin{frame}{Parte 2: Construcción de Escenarios}

\textbf{Definir dos escenarios:}

\vspace{0.3cm}

\textbf{Baseline:}
\begin{itemize}
\item Condiciones económicas normales
\item Retornos estables
\item Inflación controlada
\end{itemize}

\vspace{0.3cm}

\textbf{Adverso:}
\begin{itemize}
\item Crisis financiera
\item Caída de activos
\item Alta inflación
\item Aumento de longevidad
\end{itemize}

\vspace{0.4cm}

\textbf{Pregunta 3:}  
¿Qué hace que un escenario sea \textit{coherente}?

\end{frame}

%------------------------------------------------
\begin{frame}{Parte 3: Traducción a Shocks}

\begin{center}
\begin{tabular}{lcc}
\hline
Variable & Baseline & Adverso \\
\hline
Retorno & 5\% & -15\% (shock inicial) \\
Inflación & 2\% & 8\% \\
Longevidad & 15 años & 17 años \\
\hline
\end{tabular}
\end{center}

\vspace{0.4cm}

\textbf{Pregunta 4:}
\begin{itemize}
\item ¿Qué variable tiene mayor impacto potencial?
\item ¿Qué ocurre si el shock persiste más de un período?
\end{itemize}

\end{frame}

%------------------------------------------------
\begin{frame}{Parte 4: Modelo Previsional}

\[
W_{t+1} = W_t (1 + r_t) + c
\]

\vspace{0.3cm}

\begin{itemize}
\item $W_t$: ahorro acumulado
\item $r_t$: retorno
\item $c$: contribución
\end{itemize}

\vspace{0.4cm}

\textbf{Tarea:}
\begin{itemize}
\item Simular riqueza final bajo ambos escenarios
\item Calcular tasa de reemplazo
\end{itemize}

\end{frame}

%------------------------------------------------
\begin{frame}{Parte 5: Métricas Clave}

\textbf{Calcular:}

\begin{itemize}
\item Tasa de reemplazo esperada
\item Diferencia baseline vs adverso
\item Probabilidad de insuficiencia (extensión)
\end{itemize}

\vspace{0.4cm}

\textbf{Pregunta 5:}
\begin{itemize}
\item ¿Qué explica la caída en la tasa de reemplazo?
\item ¿Es más importante el shock inicial o la persistencia?
\end{itemize}

\end{frame}
%-----------------------------------------------
\begin{frame}[fragile]{Código R: Escenario Determinístico}
\tiny
\begin{verbatim}
#----------------------------------------
# Parámetros generales
#----------------------------------------
T <- 30              # horizonte (años)
W0 <- 10000          # ahorro inicial
c <- 1000            # contribución anual
salary <- 20000      # salario base

#----------------------------------------
# Escenario BASELINE
#----------------------------------------
r_base <- rep(0.05, T)   # retornos constantes (5%)
L_base <- 15             # años de retiro

#----------------------------------------
# Escenario ADVERSO
#----------------------------------------
r_adv <- c(-0.15, rep(0.03, T-1))  # shock inicial + recuperación lenta
L_adv <- 17                        # mayor longevidad

#----------------------------------------
# Función de simulación
#----------------------------------------
simulate_path <- function(r, W0, c, T){
  W <- W0
  for(t in 1:T){
    W <- W * (1 + r[t]) + c
  }
  return(W)
}
\end{verbatim}
   
\end{frame}
%-----------------------------------------------
\begin{frame}[fragile]{Código R: Escenario Determinístico}
\tiny
\begin{verbatim}
    
#----------------------------------------
# Simulación
#----------------------------------------
W_base <- simulate_path(r_base, W0, c, T)
W_adv  <- simulate_path(r_adv,  W0, c, T)

#----------------------------------------
# Tasa de reemplazo
#----------------------------------------
RR_base <- W_base / (L_base * salary)
RR_adv  <- W_adv  / (L_adv  * salary)

#----------------------------------------
# Resultados
#----------------------------------------
cat("Escenario Baseline - RR:", RR_base, "\n")
cat("Escenario Adverso  - RR:", RR_adv, "\n")
cat("Caída en RR:", RR_base - RR_adv, "\n")

\end{verbatim}
    
\end{frame}
%------------------------------------------------
\begin{frame}{Extensión: Simulación Monte Carlo}

\textbf{Motivación:}
\begin{itemize}
\item El futuro no es único
\item Necesitamos múltiples trayectorias
\end{itemize}

\vspace{0.4cm}

\textbf{Simular:}
\begin{itemize}
\item 1000 trayectorias de retornos
\item Distribución de la tasa de reemplazo
\end{itemize}

\vspace{0.4cm}

\textbf{Pregunta 6:}
\begin{itemize}
\item ¿Qué cambia al pasar de un escenario a una distribución?
\end{itemize}

\end{frame}

%------------------------------------------------
\begin{frame}{Métricas de Riesgo}

\begin{itemize}
\item Valor esperado
\item VaR (percentil 5\%)
\item Expected Shortfall
\item Probabilidad de insuficiencia
\end{itemize}

\vspace{0.4cm}

\textbf{Pregunta 7:}
\begin{itemize}
\item ¿Dónde está el riesgo relevante?
\item ¿Por qué la cola es crítica en pensiones?
\end{itemize}

\end{frame}
%---------------------------------------------------
\begin{frame}[fragile]{Código R: Simulación Monte Carlo}
\tiny
\begin{verbatim}
#----------------------------------------
# Simulación Monte Carlo
#----------------------------------------
set.seed(123)

N <- 5000
T <- 30
W0 <- 10000
c <- 1000
salary <- 20000

RR_sim <- numeric(N)

for(i in 1:N){  
  W <- W0  
  for(t in 1:T){
    # retornos aleatorios (media 5%, volatilidad 15%)
    r <- rnorm(1, mean = 0.05, sd = 0.15)
    W <- W * (1 + r) + c
  }  
  RR_sim[i] <- W / (15 * salary)
}


#----------------------------------------
# Métricas de riesgo
#----------------------------------------
mean_RR <- mean(RR_sim)
VaR_5   <- quantile(RR_sim, 0.05)
ES_5    <- mean(RR_sim[RR_sim <= VaR_5])
prob_insuf <- mean(RR_sim < 0.3)  # umbral de suficiencia


\end{verbatim}
    
\end{frame}
%------------------------------------------------
\begin{frame}[fragile]{Código R: Simulación Monte Carlo}
\tiny
\begin{verbatim}
    
#----------------------------------------
# Histograma
#----------------------------------------
hist(RR_sim,
     breaks = 50,
     main = "Distribución de la Tasa de Reemplazo",
     xlab = "Tasa de Reemplazo",
     border = "white")

# Línea VaR
abline(v = VaR_5, lwd = 2, lty = 2)

# Línea ES
abline(v = ES_5, lwd = 2, lty = 3)

#----------------------------------------
# Leyenda
#----------------------------------------
legend("topright",
       legend = c("VaR 5%", "Expected Shortfall"),
       lty = c(2,3),
       lwd = 2)

\end{verbatim}
    
\end{frame}

%------------------------------------------------
\begin{frame}{Comparación: Determinístico vs Monte Carlo}

\begin{itemize}
\item 1 escenario $\rightarrow$ un resultado
\item Monte Carlo $\rightarrow$ distribución completa
\end{itemize}

\vspace{0.4cm}

\textbf{Insight:}
\begin{itemize}
\item El riesgo no es un número, es una distribución
\end{itemize}

\vspace{0.4cm}

\textbf{Pregunta 8:}
\begin{itemize}
\item ¿Qué información adicional aporta Monte Carlo?
\end{itemize}

\end{frame}

%------------------------------------------------
\begin{frame}{Dinámica en Clase}

\textbf{Trabajo en grupos:}

\begin{itemize}
\item Grupo 1: escenario base
\item Grupo 2: escenario adverso moderado
\item Grupo 3: escenario adverso severo
\end{itemize}

\vspace{0.4cm}

\textbf{Cada grupo debe:}
\begin{itemize}
\item Definir narrativa económica
\item Traducir a shocks
\item Simular resultados
\item Presentar métricas clave
\end{itemize}

\vspace{0.4cm}

\textbf{Discusión:}
¿Cuál escenario representa mayor riesgo sistémico?

\end{frame}

%------------------------------------------------
\begin{frame}{Evaluación Final}

\textbf{Responder:}

\begin{itemize}
\item ¿Cuál es la diferencia entre riesgo, shock y escenario?
\item ¿Por qué un escenario debe ser coherente?
\item ¿Qué limitaciones tiene el enfoque determinístico?
\item ¿Qué aporta Monte Carlo al stress testing?
\end{itemize}

\vspace{0.5cm}

\centering
\textbf{Conclusión:}  
Los escenarios transforman riesgos en decisiones

\end{frame}

%------------------------------------------------

\end{document}